\documentclass[11pt,a4paper]{article}
\usepackage{amsmath,amssymb}
\usepackage[margin=2.5cm]{geometry}
\title{有限样本随机调度确定性等价问题}
\author{}
\date{}
\usepackage[UTF8]{ctex}
\begin{document}
\maketitle

\section*{集合与参数}
让$I$、$A$和$R$分别表示手术、麻醉师和手术室这些有限集合。让$N_s\in\mathbb{Z}_{>0}$和$\mathcal{N}=\{1,\ldots,N_s\}$分别表示样本计数和样本索引集。将$I^2_{\ne}=\{(i,j)\in I^2:i\ne j\}$和$I^3_{\ne}=\{(i,j,h)\in I^3:i,j,h\text{ pairwise distinct}\}$分别赋值。序列化的过程中使用这些不同的索引；如果存储了对角序列变量，则将其设置为零。这是根据之前的手术解释以及源PDF第7页上的明确$j\in I\setminus\{i\}$双重索引进行的。

输入参数如下：
\begin{itemize}
\item $d_i^n\geq 0$：手术的采样持续时间，为$i$，用于$(i,n)\in I\times\mathcal{N}$。
\item $c_r$代表$r\in R$，$c_a$代表$a\in A$：分别出现在手术室开启成本和麻醉师激活成本中的目标系数。
\item $c_{oa}$代表$a\in A$，$c_{or}$代表$r\in R$：分别为麻醉师和病房加班成本系数。
\item $g_a,w_a\in\{0,1\}$对于$a\in A$：常规值班和待命可用性指标。
\item $\kappa^A_{ia}\in\{0,1\}$代替$(i,a)\in I\times A$，$\kappa^R_{ir}\in\{0,1\}$代替$(i,r)\in I\times R$：分配兼容性指标。
\item $t_a^{\mathrm{start}},t_a^{\mathrm{end}}$对于$a\in A$：麻醉师班次开始和结束时间；$T^{\mathrm{end}}$：房间常规结束时间。
\item $M,M_{\mathrm{seq}},M_{\mathrm{anesth}},M_{\mathrm{room}}$: 固定有限的大$M$系数，随实例提供。
\end{itemize}
所有数值系数和集合成员都作为实例数据提供。

\section*{决策变量}
共享的第一阶段二进制变量包括：$x_{ia}$（麻醉师分配）、$y_a$（待命激活）、$z_{ir}$（房间分配）、$v_r$（房间激活）、$u_{ij}$（优先级），以及资源特定的序列变量$\alpha_{ija}$和$\beta_{ijr}$。连续变量$s_i$表示预定的开始时间。

对于每一个样本$n\in\mathcal{N}$，$S_i^n$表示实际的开始时间。分别用$O_a^{A,n}$和$O_r^{R,n}$表示源变量$\mathrm{Over}_a^n$和$\mathrm{Over}_r^n$的对应值。这些变量分别属于麻醉师和手术室加班时间两个不同的类别，无论这些变量所代表的资源标签是否重叠。

\section*{确定性等价}
所选的模型为有限样本模型，其中每个样本的最小暴露值被作为一个独立的连续变量块来对待。所有第一阶段变量在样本之间是共享的。需要对所有声明的变量进行最小化处理。
\begin{equation}
\min\quad
\sum_{r\in R}c_r v_r+\sum_{a\in A}c_a y_a
+\frac{1}{N_s}\sum_{n\in\mathcal{N}}
\left(\sum_{a\in A}c_{oa}O_a^{A,n}
+\sum_{r\in R}c_{or}O_r^{R,n}\right).
\end{equation}
术语$\sum_{r\in R}c_rv_r$仅用于表示已打开的房间，从而恢复源方程（1）。

\subsection*{第一阶段约束条件}
\begin{align}
\sum_{a\in A}x_{ia}&=1 &&\forall i\in I,\\
\sum_{r\in R}z_{ir}&=1 &&\forall i\in I,\\
z_{ir}&\leq v_r &&\forall (i,r)\in I\times R,\\
x_{ia}&\leq g_a+y_a &&\forall (i,a)\in I\times A,\\
y_a&\leq w_a &&\forall a\in A,\\
s_i&\geq t_a^{\mathrm{start}}-M(1-x_{ia}) &&\forall (i,a)\in I\times A,\\
x_{ia}&\leq\kappa^A_{ia} &&\forall (i,a)\in I\times A,\\
z_{ir}&\leq\kappa^R_{ir} &&\forall (i,r)\in I\times R.
\end{align}
\begin{align}
\alpha_{ija}&\leq u_{ij} &&\forall (i,j,a)\in I^2_{\ne}\times A,\\
\beta_{ijr}&\leq u_{ij} &&\forall (i,j,r)\in I^2_{\ne}\times R,\\
u_{ij}+u_{ji}&\leq 1 &&\forall (i,j)\in I^2_{\ne},\\
u_{ih}&\geq u_{ij}+u_{jh}-1 &&\forall (i,j,h)\in I^3_{\ne},\\
\alpha_{ija}+\alpha_{jia}&\leq x_{ia} &&\forall (i,j,a)\in I^2_{\ne}\times A,\\
\alpha_{ija}+\alpha_{jia}&\leq x_{ja} &&\forall (i,j,a)\in I^2_{\ne}\times A,\\
\alpha_{ija}+\alpha_{jia}&\geq x_{ia}+x_{ja}-1 &&\forall (i,j,a)\in I^2_{\ne}\times A,\\
\beta_{ijr}+\beta_{jir}&\leq z_{ir} &&\forall (i,j,r)\in I^2_{\ne}\times R,\\
\beta_{ijr}+\beta_{jir}&\leq z_{jr} &&\forall (i,j,r)\in I^2_{\ne}\times R,\\
\beta_{ijr}+\beta_{jir}&\geq z_{ir}+z_{jr}-1 &&\forall (i,j,r)\in I^2_{\ne}\times R.
\end{align}
\begin{align}
\alpha_{ija}&\geq x_{ia}+x_{ja}+\beta_{ijr}-2
&&\forall (i,j,r,a)\in I^2_{\ne}\times R\times A,\\
\beta_{ijr}&\geq z_{ir}+z_{jr}+\alpha_{ija}-2
&&\forall (i,j,r,a)\in I^2_{\ne}\times R\times A.
\end{align}

\subsection*{示例性约束条件}
\begin{align}
S_j^n&\geq S_i^n+d_i^n-M_{\mathrm{seq}}(1-u_{ij})
&&\forall (i,j,n)\in I^2_{\ne}\times\mathcal{N},\\
S_i^n&\geq s_i
&&\forall (i,n)\in I\times\mathcal{N},\\
O_a^{A,n}&\geq S_i^n+d_i^n-t_a^{\mathrm{end}}
-M_{\mathrm{anesth}}(1-x_{ia}+y_a)
&&\forall (i,a,n)\in I\times A\times\mathcal{N},\\
O_r^{R,n}&\geq S_i^n+d_i^n-T^{\mathrm{end}}
-M_{\mathrm{room}}(1-z_{ir})
&&\forall (i,r,n)\in I\times R\times\mathcal{N}.
\end{align}

\subsection*{变量域}
\begin{align}
x_{ia}&\in\{0,1\} &&\forall (i,a)\in I\times A,\\
y_a&\in\{0,1\} &&\forall a\in A,\\
z_{ir}&\in\{0,1\} &&\forall (i,r)\in I\times R,\\
v_r&\in\{0,1\} &&\forall r\in R,\\
u_{ij}&\in\{0,1\} &&\forall (i,j)\in I^2_{\ne},\\
\alpha_{ija}&\in\{0,1\} &&\forall (i,j,a)\in I^2_{\ne}\times A,\\
\beta_{ijr}&\in\{0,1\} &&\forall (i,j,r)\in I^2_{\ne}\times R,\\
s_i&\in\mathbb{R}_{\geq0} &&\forall i\in I,\\
S_i^n&\in\mathbb{R}_{\geq0} &&\forall (i,n)\in I\times\mathcal{N},\\
O_a^{A,n}&\in\mathbb{R}_{\geq0} &&\forall (a,n)\in A\times\mathcal{N},\\
O_r^{R,n}&\in\mathbb{R}_{\geq0} &&\forall (r,n)\in R\times\mathcal{N}.
\end{align}
\end{document}
